\subsection{Dolbeault算子及其成立条件}
在复流形上，按全纯坐标定义
\begin{align}
 \partial(f\,\dd z^I\wedge\dd\bar z^J)
 &=\sum_a\partial_a f\,\dd z^a\wedge\dd z^I\wedge\dd\bar z^J,\notag\\
 \bar\partial(f\,\dd z^I\wedge\dd\bar z^J)
 &=\sum_a\partial_{\bar a}f\,\dd\bar z^a\wedge\dd z^I\wedge\dd\bar z^J.
 \label{eq:dolbeault-local}
\end{align}
第二式如需重排为约定基，须移动$\dd\bar z^a$穿过全部$\dd z^I$，符号为$(-1)^{|I|}$。全纯换元不混合$\dd z$和$\dd\bar z$，故这两个算子就是$\dd$的$(p+1,q)$和$(p,q+1)$投影，与坐标无关。

\Needspace{14\baselineskip}
\begin{proposition}\label{prop:dolbeault-complex}
复流形上$\dd=\partial+\bar\partial$，且
\[
 \partial^2=0,\qquad\bar\partial^2=0,\qquad
 \partial\bar\partial+\bar\partial\partial=0.
\]
两算子都按总次数满足分次Leibniz规则，复共轭交换$\partial$和$\bar\partial$。
\end{proposition}
\begin{proof}
在坐标基上$\dd\dd z^j=\dd\dd\bar z^j=0$，所以只需对系数求导，得到分解。再把$\dd^2=0$分别投影到三种不同型别，得三条恒等式；或者直接用二阶偏导可交换和楔积反对称验证。分次乘法式与共轭关系也逐个坐标单项验证，因已知自然性而成为全局结论。
\end{proof}
例如$\partial\bar\partial f=\sum_{a,b}\partial_a\partial_{\bar b}f\,\dd z^a\wedge\dd\bar z^b$，而$\bar\partial\partial f$为其负号。定义
\[
 H^{p,q}_{\bar\partial}(M)=
 \frac{\ker(\bar\partial:\Omega^{p,q}\to\Omega^{p,q+1})}
 {\operatorname{im}(\bar\partial:\Omega^{p,q-1}\to\Omega^{p,q})},
\]
分母在$q=0$时为零。这里仅定义Dolbeault上同调，不预先宣称它与某个层上同调相同，亦不从第14章自动推出$H^k_{\mathrm{dR}}=\bigoplus_{p+q=k}H^{p,q}_{\bar\partial}$。

在仅有几乎复结构时，投影仍然存在，但$\dd$还可能含$(p+2,q-1)$和$(p-1,q+2)$分量。因为对一形式$\alpha$有$\dd\alpha(U,V)=U\alpha(V)-V\alpha(U)-\alpha([U,V])$，若$U,V\in T^{0,1}$而$[U,V]$出现$(1,0)$分量，$(1,0)$形式的外微分便出现$(0,2)$项。第17.4节将给出显式反例，故不能在几乎复情形不加条件地沿用本命题。

\subsection{全纯向量丛及其自然算子}
全纯向量丛是过渡矩阵全纯的复向量丛。取全纯局部标架$e$，对系数列$a$定义$\bar\partial_E(ea)=e\bar\partial a$。若$e'=eG$，全纯性$\bar\partial G=0$使定义相容，且$\bar\partial_E^2=0$。$T^{1,0}M$由全纯坐标标架及其Jacobi矩阵成为全纯丛，其对偶与外幂也一样；$K_M=\Lambda^{m,0}T^*M$称典范线丛。

反向，“光滑丛上的一个$(0,1)$联络满足平方零，便能找到全纯标架”是Koszul--Malgrange型可积性结论，本章不证明也不在构造中使用。下节Chern联络的起点始终是\emph{已经给定}的全纯丛，而不是仅给一族形式后默认为已存在全纯结构。
